\begin{definitionbox}{Definition --- Topological Space}
\begin{definition}[Topological Space]\label{def:topological-space}
Let \(X\) be a set.  A \emph{topology} on \(X\) is a collection
\(\tau \subseteq \mathcal{P}(X)\) such that:
\begin{enumerate}
  \item \(\emptyset \in \tau\) and \(X \in \tau\);
  \item if \(\{U_i\}_{i \in I}\) is any family of sets in \(\tau\), then
        \(\bigcup_{i \in I} U_i \in \tau\);
  \item if \(U_1,\ldots,U_n \in \tau\), then
        \(\bigcap_{k=1}^{n} U_k \in \tau\).
\end{enumerate}
A \emph{topological space} is a pair \((X,\tau)\), where \(X\) is a set and
\(\tau\) is a topology on \(X\).  The members of \(\tau\) are called the
\emph{open sets} of the space.
\end{definition}
\end{definitionbox}

\begin{remark*}[Standard quantified statement]
\[
\begin{aligned}
\operatorname{IsTopology}(X,\tau)
\Longleftrightarrow\;&
\tau\subseteq\mathcal{P}(X)
\land
\emptyset\in\tau
\land
X\in\tau
\\
&{}\land
\forall \mathcal{A}\subseteq\tau{:}\;
\bigcup_{U\in\mathcal{A}}U\in\tau
\\
&{}\land
\forall U_1,\ldots,U_n\in\tau{:}\;
\bigcap_{k=1}^{n}U_k\in\tau.
\end{aligned}
\]
\end{remark*}

\begin{remark*}[Predicate reading]
\[
\operatorname{IsTopology}(X,\tau)
\]
means that \(\tau\) is a topology on the carrier set \(X\).
\end{remark*}

\begin{remark*}[Interpretation]
A topology records which subsets are treated as open.  The axioms say that
openness includes the impossible and total regions, is preserved under
arbitrary unions, and is preserved under finite intersections.
\end{remark*}

\begin{remark*}[Examples]
The same set may carry different topologies.  If \(\tau_1\) and \(\tau_2\)
are different topologies on the same set \(X\), then \((X,\tau_1)\) and
\((X,\tau_2)\) are different topological spaces because they generally
declare different subsets of \(X\) to be open.
\end{remark*}

\begin{dependencies}
\begin{itemize}
  \item \hyperref[def:metric-open-set]{Definition: Metric Open Set}
\end{itemize}
\end{dependencies}

\begin{definitionbox}{Definition --- Discrete Topology}
\begin{definition}[Discrete Topology]\label{def:discrete-topology}
Let \(X\) be a set.  The \emph{discrete topology} on \(X\) is the topology
\[
  \tau_{\mathrm{disc}}=\mathcal{P}(X).
\]
Thus every subset of \(X\) is open in the topological space
\((X,\tau_{\mathrm{disc}})\).
\end{definition}
\end{definitionbox}

\begin{remark*}[Standard quantified statement]
\[
\tau_{\mathrm{disc}}=\mathcal{P}(X),
\qquad
U\text{ is open in }(X,\tau_{\mathrm{disc}})
\Longleftrightarrow
U\subseteq X.
\]
\end{remark*}

\begin{remark*}[Interpretation]
The discrete topology makes the strongest possible declaration of openness:
no subset is excluded.  It is the topology with the most open sets on the
underlying set \(X\).  Equivalently, it is the \emph{finest} topology on
\(X\): every topology on \(X\) is contained in \(\tau_{\mathrm{disc}}\).
\end{remark*}

\begin{dependencies}
\begin{itemize}
  \item \hyperref[def:topological-space]{Definition: Topological Space}
\end{itemize}
\end{dependencies}

\begin{definitionbox}{Definition --- Sierpinski Topology}
\begin{definition}[Sierpinski Topology]\label{def:sierpinski-topology}
Let \(X=\{0,1\}\).  The \emph{Sierpinski topology} on \(X\) is the topology
\[
  \tau_{\mathrm{S}}=\{\emptyset,\{1\},X\}.
\]
The topological space \((X,\tau_{\mathrm{S}})\) is called the
\emph{Sierpinski space}.
\end{definition}
\end{definitionbox}

\begin{remark*}[Standard quantified statement]
\[
U\text{ is open in }(X,\tau_{\mathrm{S}})
\Longleftrightarrow
U=\emptyset\text{ or }U=\{1\}\text{ or }U=X.
\]
\end{remark*}

\begin{remark*}[Interpretation]
The Sierpinski topology is the smallest nontrivial topology on a two-point
set.  It distinguishes the two points topologically: one singleton is open
and the other singleton is not.
\end{remark*}

\begin{dependencies}
\begin{itemize}
  \item \hyperref[def:topological-space]{Definition: Topological Space}
\end{itemize}
\end{dependencies}

\begin{definitionbox}{Definition --- Cofinite Topology}
\begin{definition}[Cofinite Topology]\label{def:cofinite-topology}
Let \(X\) be a set.  The \emph{cofinite topology} on \(X\) is the topology
\[
  \tau_{\mathrm{cof}}
  =
  \{\,U\subseteq X : X\setminus U\text{ is finite}\,\}\cup\{\emptyset\}.
\]
Thus a subset of \(X\) is open exactly when it is empty or its complement in
the ambient set is finite.
\end{definition}
\end{definitionbox}

\begin{remark*}[Standard quantified statement]
\[
U\text{ is open in }(X,\tau_{\mathrm{cof}})
\Longleftrightarrow
U=\emptyset
\text{ or }
X\setminus U\text{ is finite}.
\]
\end{remark*}

\begin{remark*}[Interpretation]
The cofinite topology treats almost all of \(X\) as open.  A nonempty open set
may omit only finitely many points, so on an infinite set any two nonempty
open sets must intersect.
\end{remark*}

\begin{dependencies}
\begin{itemize}
  \item \hyperref[def:topological-space]{Definition: Topological Space}
\end{itemize}
\end{dependencies}

\begin{definitionbox}{Definition --- Trivial Topology}
\begin{definition}[Trivial Topology]\label{def:trivial-topology}
Let \(X\) be a set.  The \emph{trivial topology} on \(X\) is the topology
\[
  \tau_{\mathrm{triv}}=\{\emptyset,X\}.
\]
Thus only \(\emptyset\) and \(X\) are open in the topological space
\((X,\tau_{\mathrm{triv}})\).
\end{definition}
\end{definitionbox}

\begin{remark*}[Standard quantified statement]
\[
\tau_{\mathrm{triv}}=\{\emptyset,X\},
\qquad
U\text{ is open in }(X,\tau_{\mathrm{triv}})
\Longleftrightarrow
U=\emptyset\text{ or }U=X.
\]
\end{remark*}

\begin{remark*}[Interpretation]
The trivial topology makes the weakest possible declaration of openness:
only the sets required by the topology axioms are open.  It is the topology
with the fewest open sets on the underlying set \(X\).  Equivalently, it is
the \emph{coarsest} topology on \(X\): it is contained in every topology on
\(X\).
\end{remark*}

\begin{dependencies}
\begin{itemize}
  \item \hyperref[def:topological-space]{Definition: Topological Space}
\end{itemize}
\end{dependencies}
